\Needspace{12\baselineskip}\SourceLeaf{315}

\begin{center}
[2.]
\end{center}

\begin{center}
\textsc{Gauss} to \textsc{Gerling}. Göttingen, 7 November 1830.
\end{center}

\indent\mbox{. . . . . . . . . .} I myself have never published anything about
the four remarkable points in a triangle. What I told
\textsc{Schumacher} about them in 1808, however, he had printed in the appendix to his
translation of \textsc{Carnot}, Géométrie de position, second
vol., which appeared in 1810. I am inclined to doubt that
the point obtained by the proposition you cite
will exhibit particularly elegant relations to the others, because
that theorem is only a very special case of a much
more general one.

For if one places on the sides of a given triangle $ABC$
three others $AB\gamma$, $\alpha BC$, $AC\beta$ (either all three
on the outside or all three on the inside of the triangle
$ABC$), each similar to a second given triangle $abc$,
in such a way that they are also similarly situated and the angles at
$\alpha$, $\beta$, $\gamma$ equal those at $a$, $b$, $c$, then
the lines $A\alpha$, $B\beta$, $C\gamma$ intersect
(extended forwards or backwards if necessary) at one and the same point
$D$, and the angles at $D$ ($ADB$, $BDC$, $CDA$) are
equal either to the angles at $\gamma$, $\alpha$, $\beta$ themselves or to their
adjacent supplementary angles, and are similarly situated.

This is evidently the most elegant construction for
\textsc{Pothenot}'s problem. In practice, however, one can be content
to construct a single point, e.\,g. $\alpha$, and to orient the plane table along
the line $A\alpha$. I hear from my son that
the latter method is known as the \textsc{Schulz-Montanus} method. \mbox{. . . . . . .}
